% Generated draft: requires mathematical and source review.
\documentclass[11pt]{article}
\usepackage[margin=25mm]{geometry}
\usepackage{fontspec}
\setmainfont{Latin Modern Roman}
\usepackage{amsmath,amssymb,mathtools,unicode-math}
\setmathfont{Latin Modern Math}
\usepackage{enumitem,needspace,xurl,hyperref}
\setlength{\parindent}{0pt}
\setlength{\parskip}{5pt}
\setlength{\emergencystretch}{3em}
\newcommand{\N}{\mathbb{N}}
\newcommand{\Z}{\mathbb{Z}}
\newcommand{\Q}{\mathbb{Q}}
\newcommand{\R}{\mathbb{R}}
\newcommand{\C}{\mathbb{C}}
\newcommand{\F}{\mathbb{F}}
\newcommand{\A}{\mathbb{A}}
\newcommand{\PP}{\mathbb{P}}
\newcommand{\E}{\mathbb{E}}
\newcommand{\eps}{\varepsilon}
\newcommand{\dd}{\,\mathrm{d}}
\begin{document}
\section*{UnsolvedMath ID: 701327}
\textbf{Source problem code:} catalogue entry
\subsection{Definitions and mathematical statement}
Let \(I\) be the \(2\times2\) identity matrix, and let \(e_{ij}\) be the \(2\times2\) matrix with entry \(1\) in position \((i,j)\) and all other entries zero. For each \(\alpha\in\C\), set
\[
A_\alpha=I+\alpha e_{12}=\begin{pmatrix}1&\alpha\\0&1\end{pmatrix},\qquad B=I+e_{21}=\begin{pmatrix}1&0\\1&1\end{pmatrix},
\]
and let \(G_\alpha=\langle A_\alpha,B\rangle\leq\mathrm{SL}_2(\C)\). Let \(F(a,b)\) denote the free group on the two formal generators \(a,b\), and let \(\phi_\alpha:F(a,b)\to G_\alpha\) be the surjective homomorphism determined by \(\phi_\alpha(a)=A_\alpha\) and \(\phi_\alpha(b)=B\). Define
\[
X=\{\alpha\in\C:\ker\phi_\alpha\neq\{1\}\}.
\]
Thus membership in \(X\) means that these particular generators satisfy a nontrivial reduced word relation.
\begin{enumerate}
\item[(a)] Is it true that
\[
\Q\cap(-4,4)\subseteq X?
\]
\item[(b)] Is it true that
\[
\Q\cap[27/7,4)\cap X\neq\varnothing?
\]
\end{enumerate}
Here rational numbers are regarded as complex numbers via the usual inclusion \(\Q\subset\C\). The registry names G. M. Bergman as proposer, while the supplied statement begins with an attribution to S. P. Farbman; their respective roles are not explained in the supplied record.
\subsection{Short English statement}
For each complex parameter, consider the two specified unipotent matrices. Let \(X\) contain the parameters for which these matrices satisfy a nontrivial group relation. The questions ask whether every rational parameter strictly between \(-4\) and \(4\) lies in \(X\), and whether at least one rational parameter in \([27/7,4)\) lies in \(X\).
\subsection{Sources}
\begin{itemize}
\item Local UnsolvedMath registry, permanent ID \(701327\), source code \texttt{catalogue entry}, title \emph{Kourovka Notebook Problem 17.25}. The supplied registry labels the problem \emph{open}; this label has not been independently verified.
\item \emph{The Kourovka Notebook (Unsolved Problems in Group Theory)}, original problem 17.25, 17th issue (2010); the supplied background also identifies the 21st issue edition (2026), arXiv:1401.0300v48. The set description names E. I. Khukhro and V. D. Mazurov as editors. These bibliographic details are supplied source data, not independently verified.
\item Attribution as supplied: G. M. Bergman, proposer, 2010; the statement itself begins with \textnormal{(S. P. Farbman)}.
\item Reference supplied in the statement: S. P. Farbman, \emph{Publ. Mat.}, \textbf{39} (1995), 379--391. No article title or link is supplied.
\item Supplied cross-references: Kourovka Notebook problem 4.41 in Archive, and problem 15.83. Their contents are not included in the supplied record.
\end{itemize}
\end{document}
